% Requires \usepackage{booktabs} in the preamble.
\subsection{SMILE Design Robustness: Perturbation Budget, Response Metric, and Surrogate Choice}
\label{sec:hull_smile_design_experiments}

Every fidelity figure reported in \S\ref{sec:hull_model_screening} comes from
one fixed SMILE design: perturbation counts of 1{,}024 (Feature) and 4{,}000
(Geo), a Wasserstein-1 response metric, and a single kernel-weighted Ridge
surrogate. That design was tuned against XGBoost and never independently
re-justified for the other six candidates -- in particular, it is not obvious
that a linear surrogate is the right choice for a model whose own
prediction surface is non-linear. This section tests three design axes
independently and in combination -- perturbation budget $N$, response metric,
and surrogate choice -- to establish which parts of the published design are
robust and which are specific to gradient-boosted trees. All experiments
reuse the same seed (42), temporal split, and $R^2>0.7$ predictive guard as
\S\ref{sec:hull_model_screening}; full results (989 rows across the three
experiments below) are in
\texttt{outputs/design\_experiments/design\_experiments\_summary.csv}.

\paragraph{Perturbation budget and response metric (both branches, Ridge
surrogate fixed).} Perturbation count was swept over $N\in\{512, 1{,}024,
2{,}048, 4{,}096\}$ (Feature) and $N\in\{2{,}000, 4{,}000, 6{,}000, 8{,}000\}$
(Geo), and the Wasserstein-1 response was compared against two LIME-style
point-wise alternatives (mean absolute difference, L1; mean squared
difference, L2), all computed from the same predict-loop per perturbation.
On the Geo branch, Wasserstein and L1 are exactly numerically identical at
every $N$ for every model -- a structural property of this branch's masking
scheme (perturbing a small subset of LSOAs never reorders the original and
perturbed prediction vectors relative to each other), not a model-specific
coincidence. Every model that clears the Geo $R^2>0.7$ guard plateaus by
$N=4{,}000$, with no gain from testing further to $8{,}000$. On the Feature
branch, gradient-boosted trees, SVR, and GP are flat and near-ceiling
($R^2\ge0.93$) across the entire grid and across all three response metrics;
MLP and DL never leave a low band ($R^2\in[0.03,0.29]$) at \emph{any} tested
$N$ or metric, and L2 is uniformly worse and less stable for both (MLP
reaches $R^2=-6.28$ at $N{=}2{,}048$ under L2). Neither more perturbations nor
a different response metric, on their own, changes the Feature-branch
picture for the neural candidates.

\paragraph{Full surrogate menu at fixed $N$ (both branches).} A
category-tailored menu -- weighted OLS, Bayesian Ridge, and the incumbent
kernel-weighted Ridge for every model; Huber and Elastic Net added for
boosting; ARD added for SVR/GP; narrow-bandwidth weighted OLS and Huber added
for MLP/DL -- plus six diagnostic controls (unweighted OLS, Ridge at both
regularisation extremes, an untuned strong Lasso, a shallow decision tree,
and a 20-repeat permuted-response null) was fit on the same mask/response
pools used for the published design. For boosting, SVR, and GP, every
candidate in the menu lands within a few points of the incumbent Ridge --
there is no surrogate that meaningfully outperforms it, and the permuted-null
control confirms the real fit is far above chance ($R^2\approx0$) for all of
these. For MLP and DL, by contrast, every linear or robust-regression
candidate stays inside the same low band the Ridge-only sweep found (MLP
best: $R^2=0.161$, Bayesian Ridge; DL best: $R^2=0.236$, Bayesian Ridge).
Bayesian Ridge is consistently the single best candidate for both neural
models, but the gain over the incumbent Ridge is marginal, not
transformative. Huber actively \emph{hurts} both neural models at every $N$
(near-zero or negative $R^2$), and the narrow-bandwidth variant is unstable,
occasionally catastrophically so ($R^2=-2.64$ for DL at $N{=}512$). Critically,
the gap between the best linear surrogate and the permuted-null baseline is
small for MLP/DL ($\sim$0.1--0.3 vs.\ a null near 0) where it is large
($\sim$0.9+) for every other model -- evidence the Feature-branch ablation
response itself carries little signal a \emph{linear} surrogate can extract,
not that the incumbent Ridge is simply a poor fit among linear alternatives.

\paragraph{Surrogate $\times$ perturbation-budget interaction, non-linear
surrogates (Feature branch, MLP/DL only).} The diagnostic-only shallow
decision tree in the menu above hinted at non-linearity (occasionally
beating every linear candidate despite being deliberately under-capacity, at
max depth 3). Random Forest and HistGradientBoosting -- fit both unweighted
and with the same kernel-distance sample weights as the linear
candidates -- were therefore tried across the full Feature $N$ grid,
extended to $N=8{,}192$ for MLP. The result overturns the "hard ceiling"
reading of the neural-network fidelity collapse for DL specifically: Random
Forest's $R^2$ climbs monotonically with $N$ (0.395 $\to$ 0.760 $\to$ 0.825
$\to$ 0.848 at $N=512,1{,}024,2{,}048,4{,}096$) -- the classic
more-perturbations-help signature that never appeared under any linear
surrogate. DL's Feature-branch ablation response is real, learnable signal;
the collapse reported in \S\ref{sec:hull_model_screening} is an artefact of
forcing a linear surrogate onto a non-linear response surface, not a
property of DL's predictions themselves. MLP shows the same qualitative
climb (0.176 $\to$ 0.261 $\to$ 0.428 $\to$ 0.474 at $N=512,1{,}024,4{,}096,
8{,}192$, unweighted Random Forest) but at a markedly lower ceiling than DL,
and with one reseed-sensitive dip at $N=2{,}048$ ($R^2=-0.747$) confirmed, by
refitting the same pool under five different train/validation/test split
seeds, to be a fold-specific artefact rather than evidence against the
trend (every other seed gives $R^2\in[0.04,0.27]$ at that same $N$).
Kernel-weighting the tree surrogates gives a small, consistent edge
(typically $+0.01$ to $+0.05$) but does not change the shape of either
model's trend. Pushing $N$ further for MLP shows visibly diminishing
returns ($+0.25$ from $N{=}512\to4{,}096$; only $+0.05$ from
$N{=}4{,}096\to8{,}192$).

\paragraph{Is the neural-network collapse architecture-general? Three further
candidates.} The MLP/DL result above raises an obvious question: is a
mask-based linear surrogate fundamentally mismatched to neural networks as a
model class, or specific to those two particular configurations? Three
further candidates were added and run through the same fidelity-only
protocol (N-sweep, full surrogate menu, surrogate$\times N$ interaction --
not the full five-metric suite of \S\ref{sec:hull_model_screening}, so these
three are not yet comparable to Table~\ref{tab:hull_model_screening} on ATT
accuracy, faithfulness, stability, or consistency): a second, deliberately
distinct feedforward ANN (wider than "dl", 256/128 vs.\ 128/64, with
LayerNormalization in place of dropout-only regularisation); a 1D-convolutional
adaptation for flat tabular input (each of the 30 features treated as a
length-30 single-channel sequence); and a lightweight per-feature
self-attention model (each feature linearly projected to a 16-dimensional
token with a learned positional embedding, one transformer-encoder block,
then pooled). All three pass the Feature-branch $R^2>0.7$ guard (0.866,
0.844, 0.814); on Geo, the ANN and transformer pass (0.721, 0.838) but the
CNN does not (0.661) -- its Geo branch was run anyway, forced past the
guard, specifically to check whether SMILE fidelity tracks predictive
quality or is independent of it.

None of the three reproduce the MLP/DL collapse. All three sit in the
$R^2\in[0.91,0.99]$ band on \emph{both} branches with the plain incumbent
Ridge surrogate at every tested $N$ -- squarely in the boosting/kernel/GP
tier, not the MLP/DL tier (Feature: ANN 0.955, CNN 0.915, Transformer 0.991
at $N{=}1{,}024$; Geo: ANN 0.943, CNN 0.943, Transformer 0.945 at
$N{=}4{,}000$). The full surrogate menu changes nothing: every candidate
clusters within a few points of Ridge for all three, exactly the pattern
already seen for boosting/kernel/GP, and unlike MLP/DL, tree-based
surrogates give no rescue because there is nothing to rescue --
HistGradientBoosting gives only a marginal edge (e.g.\ ANN Feature $0.965$ vs.\
Ridge's $0.955$ at $N{=}1{,}024$), and unweighted Random Forest actually
\emph{underperforms} Ridge on the Geo branch for all three ($R^2\in
[0.56,0.68]$ vs.\ $\approx0.94$) -- the reverse of what tree surrogates did
for MLP/DL's broken Feature branch, where they were the only thing that
worked. Fidelity is flat across the entire $N$ grid on both branches for all
three, with no plateau-then-jump signature -- these were never
under-sampled to begin with.

CNN's forced Geo result is itself informative: despite a mediocre predictive
$R^2=0.661$ that fails the guard everywhere else in this project, GeoSMILE's
fidelity there ($0.943$ at $N{=}4{,}000$) is indistinguishable from ANN's and
the transformer's gate-passing Geo branches. Surrogate fidelity measures how
well a linear model can reconstruct a black-box model's own masking
response, not how good that black-box model's predictions are -- a mediocre
predictor can still have a highly linear, easily-explained ablation
response, and this is a direct empirical case of it. It does not license
dropping the predictive guard in general (a mediocre model is still not
worth explaining for its own sake), but it confirms fidelity and predictive
quality are separate axes, not a package deal.

Taken together, these three results overturn the working explanation offered
after the original MLP/DL finding: the collapse is not a property of "neural
networks" as a class presenting a harder ablation surface to a linear
surrogate. The ANN candidate here is architecturally close to "dl" (same
dense feedforward family, deeper and wider, differing mainly in
normalization choice) and shows none of its collapse; CNN and the
transformer are structurally further from either failing model and also
show none of it. Whatever made the sklearn MLP and the original Keras "dl"
configuration collapse is specific to those two configurations -- plausibly
their comparatively narrow capacity (64--32 and 128--64--32 respectively,
against 256--128--64--32 for the ANN here) or their particular
regularisation choices, both directions this section does not attempt to
isolate further -- rather than a general obstacle to explaining deep models
with GeoSMILE.

\paragraph{Combined-model architecture check.} The same three-surrogate
comparison (Ridge, Random Forest, HistGradientBoosting), repeated for the
combined-model architecture in which a single model trained on all 32 inputs
is explained along both branches rather than two independently-trained
per-branch models (\texttt{outputs/design\_experiments/combined\_model\_
surrogate\_comparison.csv}), shows the same qualitative pattern: no
surrogate meaningfully beats Ridge for the boosting/kernel candidates, while
tree-based surrogates outperform Ridge for the neural candidates' Feature
branch. This rules out "the two-branch architecture specifically causes the
neural-network collapse" as an explanation -- the effect is a property of
the model class and the linear surrogate, not of which architecture explains
it.

\paragraph{Correction to \S\ref{sec:hull_model_screening}.} While tracing
these results, MLP's Geo-branch predictive $R^2$ was confirmed at $0.444$
(\texttt{outputs/eval\_by\_model/predictive\_performance.csv}), below the
$0.7$ guard applied everywhere else in this project. The published values
this section is correcting had been computed before that guard was
introduced and were never re-screened against it; Table
\ref{tab:hull_model_screening} and the accompanying prose have been updated
to exclude MLP's Geo branch on the same terms as SVR's, and a stray
comparison citing MLP's now-invalid Geo consistency figure has been removed.

\paragraph{Implications.} These experiments justify keeping the published
design (kernel-weighted Ridge, $N=1{,}024$/$4{,}000$) for every model in the
boosting and kernel/GP categories: no tested surrogate, response metric, or
budget improves on it, and the permuted-null control confirms the fit
reflects real signal rather than surrogate overfitting. For MLP and DL, the
published Feature-branch exclusion remains the correct call under a linear
surrogate at any tested budget -- but the surrogate, not the budget or
response metric, was the actual bottleneck: a tree-based surrogate very
nearly closes DL's Feature-branch fidelity gap ($R^2=0.848$ at $N{=}4{,}096$,
against $\ge0.93$ for every boosting/kernel model) and substantially narrows
MLP's ($R^2=0.474$ at $N{=}8{,}192$, against a linear-surrogate ceiling of
$0.161$). Adopting a tree surrogate for these two models is not, however, a
drop-in fix: GeoSMILE's reported per-feature attribution is the Ridge
surrogate's linear coefficient, which Random Forest and HistGradientBoosting
have no equivalent of, so realising this fidelity gain as an actual
attribution ranking would require a separate importance method (e.g.\
permutation or SHAP importance computed on the tree surrogate) layered on
top -- a natural direction for follow-up work, not pursued here. Finally,
the ANN/CNN/transformer results reframe what "the neural-network problem"
actually is: it is not that GeoSMILE struggles with deep or non-linear
predictors in general -- three further architectures, one of them
attention-based, show fidelity indistinguishable from gradient-boosted
trees -- but that two specific, comparatively narrow configurations (the
screened MLP and DL) happen to produce a masking response a linear
surrogate cannot capture. This matters for how the published exclusion of
MLP/DL should be described: not as a limitation of GeoSMILE against neural
networks, but as a property of those two trained models.
